\documentclass[output=paper]{langscibook}
\ChapterDOI{10.5281/zenodo.22880228}
\title{The Teop language documentation and the content of this volume}
\author{Ulrike Mosel}
\abstract{\noabstract}
\begin{document}
\maketitle\label{chap:2}
\section{The historical background of the Teop language documentation}
The Teop language documentation project started with a fieldwork methods course which the native speaker Ruth Saovana Spriggs and I gave at the Australian National University in 1992. In 1994 we got a grant from the Australian Research Council that enabled Saovana to travel home and do audio-recorded interviews about the culture of her people and their personal histories. When I moved to Kiel in Germany in 1995, Saovana and I continued our work per email and in person when she visited me in Kiel for three months in 1995 and two months in 1999.

From 2000 to 2007 our project was generously funded by the Dokumentation-Bedrohter-Sprachen Programm (DobeS) of the Volkswagen Foundation, so that we could do further fieldwork, pay the Teop team members in Bougainville for their full time or part time work, and employ research assistants in Kiel and Canberra.\footnote{See \url{https://dobes.mpi.nl/projects/teop/team/}.} Saovana was a team member until March 2005 when she left the project to concentrate on her PhD thesis. In 2009 the project was funded for 16 months by the German Research Foundation, and afterwards by private resources. In 2014, I visited Bougainville for the last time. 

The organization of our work in Bougainville was considerably influenced by the so-called Bougainville Conflict, a civil war which had its roots in the people’s sufferings from colonialism and started in 1988. The most important dates of the island’s history are listed in \tabref{tab:history}.

\begin{table}
    \caption[Some dates of the history of the Bougainville Island]{Some dates of the history of the Bougainville Island. See \textcite{Sack:2015hy}, \textcite{Nelson:2015fl}, \textcite{Mahaka:2014oc}, \textcite{Parker:2010uo}, and \textcite{Purupuru:2014la}.}
    \label{tab:history}
    \begin{tabularx}{\textwidth}{lQ}
        \lsptoprule
            \textsc{date}&\textsc{description}\\
        \midrule
            1885--1914 & German colonial rule\\
            1901--1902 & First mission of the Catholic Marists\\
            1914--1920 & Australian military administration\\
            1920--1942 & Australian administration granted by the League of Nations\\
            1920 & Start of the Methodist mission\\
            1924 & Start of the Seventh-day Adventist mission\\
            1942--1945 & 2nd World War\\
            1942 & Invasion by the Japanese army\\
            1943 & Invasion by the Allied Forces\\
            1945--1975 & Australian administration\\
            1972 & Beginning of mining at the Paguna copper and goal mine (Panguna Mine)\\
            1975 & Independence of Papua New Guinea (PNG)\\
            1988 & Begin of the Bougainville Conflict between the landowners of Panguna and the Bougainville Copper Ltd, which escalated into an armed conflict \\
            1990--1997 & Absence of an effective government body, development of anarchy\\
            1997 & Truce, deployment of peace monitoring groups from New Zealand\\
            1998 & Permanent ceasefire, peace monitoring groups from Australia\\
            2001 & Peace agreement\\
        \lspbottomrule
    \end{tabularx}
\end{table}

The Bougainville Conflict was a civil war that arose between the Papua New Guinea Security Forces (PNG SF) and Bougainvillean people who formed the Bougainville Revolutionary Army (BRA) to fight against the disastrous environmental effects of the Panguna Mine, the exploitation of the indigenous landowners, and the immigration of 15,000--20,000 non-Bougainvillean workers.\footnote{\textcite{Regan:2001ky}} When in 1989 the mine closed and in 1990 the PNGDF (Papua New Guinea Defense Force) imposed a blockade on Bougainville Island, all immigrants left the island and all kinds of trade were stopped, so that ``subsistence gardening became the sole means of survival'' (\cite[4]{Havini:2004tt}). Bougainvillean people who had lived and worked in other provinces had to return to their villages in Bougainville. The Teop people were not directly involved in the battles, but they suffered from the blockade, unjust treatment by both parties, and attacks by criminal gangs. After the peace agreement in 2001, this insecure and economically demanding situation continued for many years. 

\section{The collaborative fieldwork}
\largerpage[-1]
In 2000 Saovana and I travelled for a short visit to Bougainville where she introduced me to some Teop leaders and to Ruth Sima Rigamu who invited me to stay with her in the village Hiovabon during future fieldwork visits. After this visit, Saovana went to Bougainville by herself in 2001 and 2003 to do more interviews, whereas I did twelve fieldwork trips by myself between 2003 and 2014, and one in 2006 together with the student assistant Yvonne Thiesen. Since Saovana and I did not do fieldwork together in Bougainville, this section only describes the work of the team I worked with in Hiovabon.

The social and economic consequences of the Bougainville Conflict had a considerable impact on the organization of our fieldwork and the subsequent processing of the data collected:
\begin{enumerate}
    \item Because of the lack of any trade stores in the Teop area and the need for subsistence gardening and fishing, the average duration of our fieldwork was only four weeks.
    \item For security reasons, I was not allowed to bring a computer with me until 2008, so that all transcriptions of recordings, translations and fieldnotes had to be written by hand, then typed up and printed in Kiel, reviewed and corrected by hand during the next fieldwork, brought back to Kiel to be incorporated in the Teop Language Corpus as ELAN\footnote{ELAN is an annotation tool for audio and video recordings developed by the Max-Planck Institute für Psycholinguistics (\url{https://archive.mpi.nl/tla/elan}), see also \textcite[102-105]{Barth:2022kc}.} and PDF files, and in the Toolbox\footnote{Toolbox is a data management and analysis tool for field linguists that has been developed by SIL (\url{https://software.sil.org/toolbox/}), see also \textcite[177]{Meakins:2018bj}. The Teop Toolbox lexical data base has been organized and cared for by Marcia Schwartz from 2000 to 2007, whereas Jessika Reinig and Yvonne Thiesen worked as research assistants on the annotation of the Teop texts in ELAN from 2000--2003 and 2000--2007, respectively. The Teop phonology was analyzed and described by Alexander Radtke (\citeyear{Radtke:2005kw}).} lexical data base, and eventually brought back or sent to the Teop people as books printed in Germany (\cite{Mosel:2007qo}; \cite{Mahaka:2010kk}).
    \item Because of the Teops’ bad experiences with foreigners, neither I nor the Teop team members could ask the speakers to sign ``a participant consent form'' (\cite[37--38]{Meakins:2018bj}). I also refrained from asking personal questions about clan membership, social status, or education, so that the metadata of the Teop Language Corpus (\textsc{tlc}) documents do not conform to the ideals of digital archives (\cite[76]{Meakins:2018bj}).
\end{enumerate}

During my first two field trips in 2003 my host Ruth Sima Rigamu and the teacher and graphic artist Neville Vitahi introduced me to seven Teop villages and helped me to explain and discuss our documentation  project. Simultaneously they organized our local research team, so that we could do the first recordings and transcriptions of spontaneously spoken folktales, stories about the Teop history, and descriptions of traditional customs and practices. By the end of our second fieldwork in 2003, it became clear that the transcriptions needed to be supplemented by edited versions, because the transcriptions are neither suitable as reading materials for school children nor a pleasure to read for any readers other than linguists (\cite{Mosel:2004zp,Mosel:2024yp}). But for linguists the edited texts and their oral originals show what kind of lexical and grammatical changes the editors made when they transformed oral texts into written ones and thus created a new register in a language which previously lacked written folk tales, historical narratives, and descriptions of the environment, artefacts and crafts (\cite{Mosel:2015bs}).

\begin{figure}
    \includegraphics[width=0.9\linewidth]{figures/booklaunch.jpg}
    \caption{Launch of the school book on house building}
    \label{fig:booklaunch}
\end{figure}

In March 2004, the Teop team -- Enoch Horai Magum, Joice Maion, Naphtali Maion, Ruth Sima Rigamu and Neville Vitahi -- decided to record more folktales, and transcribe and edit them for a primary school reader, but also to continue the recording of texts about the history and the traditional culture of the Teop people, because elderly people were keen to transmit their knowledge to future generations. From 2005 onwards the team members felt confident in writing Teop texts without doing recordings and transcriptions. The edited and written texts were proof-read by the teachers Jubilee Kamai and Ondria Tavagaga. In the following years we published a school book with 40 folktales (\cite{Mosel:2007qo}), a small book on house building (\cite{Mahaka:2010kk}), a bilingual history book (\cite{Mosel:2014kg}), and much later the Multifunctional Teop-English Dictionary, abbr. \textsc{mted}, in which short encyclopedic texts were included as ``examples'' with translations, see \textcite{Mosel:2019bt}.
The Teop texts of this volume are of three types:
\begin{enumerate}
    \item originally oral texts whose speaker figures as the first author, whereas the editor figures as the second author;
    \item texts that were written and edited in team work; the first author is the one who wrote the first version, whereas the second author did most of the editorial work;
    \item texts that have been written and edited by the same author.
\end{enumerate}
The translations were always done in teamwork. They were either done by the Teop author or the editor, and then explained to me, or I did the translation with the help of the author, the editor or my host Ruth Sima Rigamu, but as the final editor of this volume I carry the responsibility for all translations.

From a corpus linguistic perspective, the Teop Language Corpus (\textsc{tlc}), and in particular this text collection is a so-called opportunistic corpus without a specified sampling frame (\cite[250]{Mosel:2018fb}).
\section{The content and the structure of this volume}
The Teop texts and their English translations are presented in two forms. Part \ref{part:2} to Part \ref{part:5} present the chapters on plants, house building, canoes and fishing techniques as bilingual texts in two columns with photographs and drawings. Each of these four parts is introduced by a short chapter which summarizes its content and gives some background information with references to the anthropological literature and to related dictionary entries in the \textsc{mted}.

Part \ref{part:6} contains the same texts and their translations, but in addition provides a morphological word-by-word segmentation of the Teop text with a linguistic glossing of each segment. If the free English translation is markedly different from its Teop source with respect to its syntax and phraseology, a literal translation is added in footnotes, so that the Teop text, its morphological segmentation, glossing, and free translation are presented in four lines, see below:

\medskip

\begin{tabular}{llll}
    Three words of line (\ref{chap:5}:\ref{ex:text01-148})&&&\\
    Teop text & \multicolumn{3}{l}{\textit{teo vuri nahu}}\\
    Word division & te=o & vuri & nahu\\
    Glossing & \textsc{prep}=\textsc{art}3.\textsc{sg} & time & cook\\
    Free translation & \multicolumn{3}{l}{`when they cook them'\footnotemark}\\
    %Literal translation & \multicolumn{3}{l}{lit. `at the time (of) cook(ing)'}\\
\end{tabular}\footnotetext{lit. `at the time (of) cook(ing)'}

\medskip

If you have the PDF-version of this book and find a quotation like ``see (\ref{chap:5}:\ref{ex:text01-148})'', you can simply click this reference and then find it in its context. The first number links to the chapter in Parts \ref{part:2}--\ref{part:5}, while the second number links to the glossed text in Part \ref{part:6}.
\section{Bibliographic data of the text collection}
For each chapter of this volume, Table \ref{tab:texts} below lists the labels of the text in the Teop Language Corpus (\textsc{tlc}) and the examples in the Multifunctional Teop-English Dictionary (\textsc{mted}), two dates, and a word count for each text. If the chapter represents the edited version (E) of an oral text, the first date refers to the recording (wav), and in the case of a written text (W) it refers to the production of the manuscript (ms). The second date refers to the production of the ELAN text in both cases.

Most chapters consist of a text that is the same as the one in the \textsc{tlc}. The exceptions are Chapters \ref{chap:13}, \ref{chap:14}, and \ref{chap:23} which are excerpts from descriptions that Ruth Sima Rigamu wrote for the \textsc{mted}.

\begin{table}
    \caption{The texts in this collection.} \label{tab:texts}
\small
\begin{tabularx}{\textwidth}{rQQllr}
   
    \lsptoprule
    \textsc{chap}&\textsc{title}&\textsc{authors}&\textsc{label}&\textsc{dates}&\textsc{words}\\
    \midrule
%     \endfirsthead
%     \multicolumn{6}{c}%
%     {\tablename\ \thetable{} -- continued from previous page}\\
%     \lsptoprule
%         \textsc{chap}&\textsc{title}&\textsc{authors}&\textsc{label}&\textsc{dates}&\textsc{words}\\
% 	    \midrule
%     \endhead
%     \hline \multicolumn{6}{r}{{Continued on next page}}\\
%     \endfoot
%     \endlastfoot
            \multicolumn{6}{l}{\textsc{part ii} -- Plants used for artefacts}\\
            \midrule
            \ref{chap:5}&\textit{O overe} -- The coconut palm &Joyce Maion &Joy\_17W&2009/2011&1,552\\
            \ref{chap:6}&\textit{O atovo} -- The sago palm&Ruth Sima Rigamu and Joyce Maion&Sii\_Joy\_03W&2009/2012&491\\
            \ref{chap:7}&\textit{O sinivi} -- The canoe tree (Blackboard Tree)&Joyce Maion&Joy\_24W&2011/2012&190\\
            \ref{chap:8}&\textit{O baneo visoas} -- The young breadfruit tree&Ruth Sima Rigamu and Joyce Maion&Sii\_Joy\_05W&2012/2013&202\\
            \ref{chap:9}&\textit{O poosusu} -- The sea poison tree&Joyce Maion&Joy\_20W&2011/2012&86\\
            \ref{chap:10}&\textit{Koverau, kururu, koopu} -- Three kinds of bamboo&Ruth Sima Rigamu and Joyce Maion&Sii\_40W&2012/2012&392\\
            \ref{chap:11}&\textit{Kara naono} -- Vines, aerial roots, fibres and strings&Ruth Sima Rigamu&Sii\_24W&2009/2011&747\\
            \ref{chap:12}&\textit{Kara roge} -- Hibiscus fibres for making strings&Joyce Maion&Joy\_18W&2009/2011&127\\
            \ref{chap:13}&\textit{Karaa} -- Pandanus&Sima Rigamu&Sii\_17W\_191-211&2009/2009&225\\
            \ref{chap:14}&\textit{Taovu} -- The giant pandanus and the hood&Ruth Sima Rigamu and Joyce Maion&Sii\_41W\_31-44&2012/2014&141\\
            \midrule
            \end{tabularx}
            \end{table}

            \begin{table}
            \small
            \begin{tabularx}{\textwidth}{rQQllr}
            \midrule
            \textsc{chap}&\textsc{title}&\textsc{authors}&\textsc{label}&\textsc{dates}&\textsc{words}\\
            \midrule
            \multicolumn{6}{l}{\textsc{part iii} -- Building houses from plant materials}\\
            \midrule
            \ref{chap:16}&\textit{A sinoo na inu} -- The frame of the house&Enoch H. Magum&Eno\_05W&2005/2008&393\\
            \ref{chap:17}&\textit{A kiu vaa teo atovo} -- The work on the sago palm thatch&Enoch H. Magum&Eno\_02W&2005/2008&705\\
            \ref{chap:18}&\textit{O kovuaka} -- The top thatch sheet&Enoch H. Magum&Eno\_07W&2005/2008&147\\
            \ref{chap:19}&\textit{A takinima} -- The woven bamboo wall&Enoch H. Magum&Eno\_03W&2005/2008&445\\
            \ref{chap:20}&\textit{O paranda} -- The floor&Enoch H. Magum&Eno\_04W&2005/2008&380\\
            \midrule
            \multicolumn{6}{l}{\textsc{part iv} -- Dugout canoes with outriggers}\\
            \midrule
            \ref{chap:22}&\textit{O sinivi suuna} -- The indigenous canoe&Mark Tahi and Enoch H. Magum&Tah\_02E(Eno)&2003/2006&357\\
            \ref{chap:23}&\textit{O aga} -- The outrigger float&Sima Rigamu&Sii\_17W; tier 031-044&2008/2009&103\\
            \ref{chap:24}&\textit{O sinivi potee} -- The imported canoe called \textit{potee}&Paul Morekevan and Enoch H. Magum&Mor\_04E(Eno)&2003/2009&612\\
            \midrule
            \end{tabularx}
            \end{table}

            \begin{table}
            \small
            \begin{tabularx}{\textwidth}{rQQll@{}r}
            \midrule
            \textsc{chap}&\textsc{title}&\textsc{authors}&\textsc{label}&\textsc{dates}&\textsc{words}\\
            \midrule
            \multicolumn{6}{l}{\textsc{part v} -- Fishing techniques}\\
            \midrule
            \ref{chap:26}&\textit{A goroto} -- The turtle&Naphtali Maion and Ondria Tavagaga&Naph\_01W&2008/2009&579\\
            \ref{chap:27}&\textit{Kave goroto} -- The turtle net&Mark Tahi and Enoch H. Magum&Eno\_08W&2003/2009&477\\
            \ref{chap:28}&\textit{A vaahoa kave goroto} -- The launching of the turtle net&Enoch H. Magum&Eno\_09W&2008/2009&213\\
            \ref{chap:29}&\textit{Kakavei} -- The bamboo frame net&Naphtali Maion&Naph\_03W&2008/2009&624\\
            \ref{chap:30}&\textit{Varigomo} -- A kind of net fishing&Enoch H. Magum&Eno\_12W&2006/2009&295\\
            \ref{chap:31}&\textit{Seige} -- Fishing with a net on two mangrove poles&Enoch H. Magum&Eno\_14W&2006/2009&241\\
            \ref{chap:32}&\textit{Tanaa} -- Fishing in a group during low tide&Helen K, Magum and Shalom Magum&Hel\_08E(Sha)&2009/2010&272\\
            \ref{chap:33}&\textit{Vavaaiku} -- Angling with a swinging bait&Enoch H. Magum and Shalom Magum&Eno\_15W&2009/2012&204\\
            \ref{chap:34}&\textit{Vavaatoa} -- Kite fishing&Naphtali Maion&Naph\_02W&2008/2009&238\\
            \ref{chap:35}&\textit{Vagana haopi} -- Catching wrasses by hand&Helen K, Magum and Shalom Magum&Hel\_09E(Sha)&2009/2010&195\\
            \ref{chap:36}&\textit{Paka suva} -- Fish poisoning leaves&Helen K, Magum, Shalom Magum and Joyce Maion&Hel\_14E(Sha)&2009/2011&214\\
            \ref{chap:37}&\textit{Ana suva} -- Fish poisoning roots&Helen K, Magum, Shalom Magum and Joyce Maion&Hel\_15E(Sha)&2009/2011&201\\
            \ref{chap:38}&\textit{Suvasuva kokobus} -- Fish poisoning with sea cucumbers&Helen K, Magum, Shalom Magum and Joyce Maion&Hel\_16E(Sha)&2009/2011&161\\
            \ref{chap:39}&\textit{Hovuhovu iobo} -- Digging sea-sausages&Helen K. Magum and Shalom Magum&Hel\_07E(Sha)&2009/2010&303\\
            \midrule
            \multicolumn{5}{l}{\textbf{Total}}&\textbf{11,512}\\
        \lspbottomrule
    \end{tabularx}
\end{table}
\end{document}